% ----- DOHA 1 -----
\subsection*{\deva{प्रथम दोहा} (First Doha)}
\addcontentsline{toc}{subsection}{\deva{प्रथम दोहा} (First Doha) — \deva{श्रीगुरु चरन...}}

\begin{minipage}[t]{0.55\textwidth}
\vspace{0pt}

{\large \linespread{1.5}\selectfont
\makebox[\linewidth][l]{\deva{श्रीगुरु चरन सरोज रज}}\\
\makebox[\linewidth][c]{\deva{निज मनु मुकुरु सुधारि।}}\\
\makebox[\linewidth][l]{\deva{बरनउँ रघुबर बिमल जसु}}\\
\makebox[\linewidth][c]{\deva{जो दायकु फल चारि॥}}
\par}

\vspace{0.5em}

{ \linespread{1.5}\selectfont
\makebox[\linewidth][l]{\telu{శ్రీగురు చరన సరోజ రజ }}\\
\makebox[\linewidth][c]{\telu{నిజ మను ముకురు సుధారి।}}\\
\makebox[\linewidth][l]{\telu{బరనఉఁ రఘుబర బిమల జసు}}\\
\makebox[\linewidth][c]{\telu{జో దాయకు ఫల చారి॥}}
\par}

\vspace{0.5em}

{ \linespread{1.3}\selectfont
\textit{śrīguru carana saroja raja, nija manu mukuru sudhāri |}\\
\textit{baranaũ raghubara bimala jasu, jo dāyaku phala cāri ||}
\par}
\end{minipage}%
\hfill
\begin{minipage}[t]{0.42\textwidth}
\vspace{0pt}
\begin{tcolorbox}[anchorbox]
\textbf{The Frame of Mind & Dedication:} Why begin with \textit{Guru Vandana} (respecting the teacher)? Because actively thinking about purifying one's mind (\textit{Mukuru Sudhari}) immediately puts the student in a positive, receptive frame of mind to learn. Hanuman perfectly modeled this respect for education: to learn the scriptures from Surya (the Sun), he flew backward all day across the blazing sky just to keep his eyes respectfully on his teacher. It shows supreme dedication to learning.
\vspace{0.5em}\newline Jambavan acts as a true mentor, wiping away self-doubt by reminding us of our immense, hidden potential.\newline\null\hfill\href{https://www.valmiki.iitk.ac.in/sloka?field_kanda_tid=4&language=dv&field_sarga_value=66}{\textit{Kishkindha Kanda, 66:31-35}}
\end{tcolorbox}
\end{minipage}

\vspace{0.5em}
\subsubsection*{\textbf{Visual Rhythm Map:}}
\begin{table}[h!]
\centering
\renewcommand{\arraystretch}{1.2}
\begin{tabularx}{\textwidth}{|>{\centering\arraybackslash\hsize=1.3333\hsize}X|>{\centering\arraybackslash\hsize=1.0000\hsize}X|>{\centering\arraybackslash\hsize=1.3333\hsize}X|>{\centering\arraybackslash\hsize=0.6667\hsize}X|>{\centering\arraybackslash\hsize=0.6667\hsize}X|>{\centering\arraybackslash\hsize=0.6667\hsize}X|>{\centering\arraybackslash\hsize=1.0000\hsize}X|>{\centering\arraybackslash\hsize=1.3333\hsize}X|}
\hline
\textbf{\guru{श्री}\laghu{गु}\laghu{रु}} \par (4) & \textbf{\laghu{च}\laghu{र}\laghu{न}} \par (3) & \textbf{\laghu{स}\guru{रो}\laghu{ज}} \par (4) & \textbf{\laghu{र}\laghu{ज}} \par (2) & \textbf{\laghu{नि}\laghu{ज}} \par (2) & \textbf{\laghu{म}\laghu{नु}} \par (2) & \textbf{\laghu{मु}\laghu{कु}\laghu{रु}} \par (3) & \textbf{\laghu{सु}\guru{धा}\laghu{रि}} \par (4) \\
\hline
\end{tabularx}
\vspace{-1.5em}
\end{table}

\begin{table}[h!]
\centering
\renewcommand{\arraystretch}{1.2}
\begin{tabularx}{\textwidth}{|>{\centering\arraybackslash\hsize=1.3333\hsize}X|>{\centering\arraybackslash\hsize=1.3333\hsize}X|>{\centering\arraybackslash\hsize=1.0000\hsize}X|>{\centering\arraybackslash\hsize=0.6667\hsize}X|>{\centering\arraybackslash\hsize=0.6667\hsize}X|>{\centering\arraybackslash\hsize=1.3333\hsize}X|>{\centering\arraybackslash\hsize=0.6667\hsize}X|>{\centering\arraybackslash\hsize=1.0000\hsize}X|}
\hline
\textbf{\laghu{ब}\laghu{र}\laghu{न}\laghu{उँ}} \par (4) & \textbf{\laghu{र}\laghu{घु}\laghu{ब}\laghu{र}} \par (4) & \textbf{\laghu{बि}\laghu{म}\laghu{ल}} \par (3) & \textbf{\laghu{ज}\laghu{सु}} \par (2) & \textbf{\guru{जो}} \par (2) & \textbf{\guru{दा}\laghu{य}\laghu{कु}} \par (4) & \textbf{\laghu{फ}\laghu{ल}} \par (2) & \textbf{\guru{चा}\laghu{रि}} \par (3) \\
\hline
\end{tabularx}
\vspace{-1.5em}
\end{table}

\vspace{0.5em}
\textbf{ \deva{सरल-संस्कृतम्}:}\\
 \deva{श्रीगुरोः चरण-कमल-धूल्या स्व-मनः-रूपी-दर्पणं स्वच्छं कृत्वा}\\
 \deva{अहं श्रीरघुनाथस्य निर्मलं यशः वर्णयामि, यत् यशः चत्वारि फलानि (धर्म-अर्थ-काम-मोक्षान्) ददाति।}\\

Having cleansed the mirror of my mind with the dust of my Guru's lotus feet\\
I narrate the pure glory of Lord Rama, which bestows the four fruits of life.


\vspace{1em}
\newpage
\textbf{\deva{प्रतिपदार्थः} (Meaning of each word):}
\begin{table}[h!]
\centering
\renewcommand{\arraystretch}{1.2}
\begin{tabularx}{\textwidth}{@{} l l X X @{}}
\toprule
\textbf{Awadhi} & \textbf{Sanskrit} & \textbf{English} & \textbf{Grammar \& Notes} \\
\midrule
\deva{श्रीगुरु} / \telu{శ్రీగురు} / \textit{śrīguru}  &  \deva{श्रीगुरोः}  &  Venerable Guru's  & Genitive (-ḥ) \\
\deva{चरन} / \telu{చరన} / \textit{carana}  &  \deva{चरण}  &  Feet  & \\ 
\deva{सरोज} / \telu{సరోజ} / \textit{saroja}  &  \deva{कमल}  &  Lotus  &   \\
\deva{रज} / \telu{రజ} / \textit{raja}  &  \deva{धूलिः / रजः}  &  Dust  &   \\
\deva{निज} / \telu{నిజ} / \textit{nija}  &  \deva{स्व / आत्मीय}  &  Own  &   \\
\deva{मनु} / \telu{మను} / \textit{manu}  &  \deva{मनः}  &  Mind  &   \\
\deva{मुकुरु} \gotchaicon / \telu{ముకురు} / \textit{mukuru}  &  \deva{दर्पणम्}  &  Mirror  & Nominative -u ending \\
\deva{सुधारि} \gotchaicon / \telu{సుధారి} / \textit{sudhāri}  &  \deva{स्वच्छं कृत्वा (संशोध्य)}  &  Having cleaned/purified  &   \\
\deva{बरनउँ} / \telu{బరనఉఁ} / \textit{baranaũ}  &  \deva{वर्णयामि}  &  I describe / narrate  & \\ 
\deva{रघुबर} / \telu{రఘుబర} / \textit{raghubara}  &  \deva{श्रीरघुनाथस्य}  &  Lord Rama's  & \\ 
\deva{बिमल} / \telu{బిమల} / \textit{bimala}  &  \deva{निर्मलम्}  &  Pure  & \\ 
\deva{जसु} \gotchaicon / \telu{జసు} / \textit{jasu}  &  \deva{यशः}  &  Glory  & \\ 
\deva{जो} / \telu{జో} / \textit{jo}  &  \deva{यत्}  &  Which  &   \\
\deva{दायकु} \gotchaicon / \telu{దాయకు} / \textit{dāyaku}  &  \deva{दायकम् (फलप्रदम्)}  &  Bestower / Giver  & Nominative -u ending \\
\deva{फल चारि} \glossaryicon / \telu{ఫల చారి} / \textit{phala cāri}  &  \deva{चत्वारि फलानि}  &  Four fruits  &   \\
\bottomrule
\end{tabularx}
\end{table}

\begin{tcolorbox}[gotchabox]
\begin{itemize}
    \item \textbf{The -u Ending (Nominative Case):} The short 'u' on \deva{मुकुरु} (Mukuru), \deva{जसु} (Jasu), and \deva{दायकु} (Dayaku) doesn't mean it's a verb---it is the classical Apabhramsha nominative case marker derived from the Sanskrit '-aḥ' ending.
    \item \textbf{Meter Alert:} Keep the 'i' short in \deva{सुधारि} (Sudhaari) and \deva{चारि} (Chaari) to maintain the exact 11-beat math of the Doha's second half! Both end in the classic Guru-Laghu (2-1) drop.
\end{itemize}
\end{tcolorbox}

\begin{tcolorbox}[poetrybox]
Notice the beautiful structural pairing of \deva{सुधारि} (\textit{sudhāri}) and \deva{चारि} (\textit{cāri}) at the end of the half-verses, establishing the melodic pace of the Doha meter.
\end{tcolorbox}


\newpage
